%Not more than 1 page, single-spaced and written in narrative format. Should include a summary of previous education, including institutions attended, and research experiences. 

\begin{singlespace}
Originally from Venezuela, I started my undergraduate studies at the University of the Andes in Merida, Venezuela in 2004 in Geological Engineering. Geological Engineering was one of the most sought after Engineering disciplines due to the country's strong petroleum industry, therefore transferring to a different discipline, if necessary, would be simpler than the other way around. Sure enough, after learning about electromagnetism in Physics II, I transferred to Mechanical Engineering. 

Halfway through my undergraduate studies, the opportunity arose to join my family in Minnesota and finish my Bachelor's degree at the University of Minnesota. While pursuing the engaging coursework in Mechanical Engineering, I had the opportunity to work in two separate research labs. In the Hydraulics Laboratory in the Dept. of Bioproducts and Biosystems Engineering, I was in charge of the research and development of a subterranean fluxmeter, a device that measures subterranean fluid flow while being minimally invasive. The fluxmeter was capable of sensing flow values comparable to those recorded by commercial rain gauges. Follwing this, I carried out cell biostabilization research in Prof. Aksan's laboratory using fluorescent microscopy to evaluate human foreskin fibroblast cell viability under a number of adversarial osmotic and temperature conditions, additionally, I created a testing device capable of suspending cells as small as 6 $\mu m $ in diameter in flow to assist in cell viability research. 

I was admitted into the Mechanical Engineering Department's Ph.D. program following my graduation from the undergraduate program in Fall of 2009 and began working in Prof. Papanikolopoulos's Center for Distributed Robotics (CDR) in the summer of 2010. My time at the CDR has been one of great learning across a wide variety of fields, from acquiring depth in Mechanical Engineering topics, to expanding across to topics typically found in Electrical and Computer Engineering, such as embedded system design and programming, development of custom embedded operating systems, integration of a micro Geiger-Muller sensor, as well as design of inexpensive plant growth chambers for genetics research on Arabidopsis thaliana. 

Research in the CDR includes the work on an amphibious robotic platform, the Aquapod. This team effort created a small and inexpensive robot for collecting data from litoral environments where the amphibious nature would permit data collection opportunities not present in other robots. Further research on this platform has led to the screwdrive drivetrain, as well as the Omnibot. The screwdrive locomotion method aims to enhance the locomotive capabilities of the existing Aquapod platform. The Omnibot platform is a sibling platform to the Aquapod yet capable of flight. This platform is the first of its kind and poses a number of technical challenges, however, the benefits in terms of aerial surveillance, loitering, and deployment capabilities cannot be understated.  

%Research Experience and Relevant Coursework
%enter for Distributed Robotics (Dept. of Computer Science) Fall 2010 – Current
%orked on robot design and sensor integration for multiple robots. Some of the projects are expansion of
%he Adelopod drivetrain to include tank treads, mechanical design, sensor and component specification, and
%ystem testing for the Aquapod robotic platform.
%iostabilization Laboratory (Dept. of Mech. Eng.) Spring 2009 – Fall 2009
%tudied human foreskin fibroblast cell viability when subjected to varying osmotic pressures and
%emperatures using fluorescent microscopy and Fourier Transform Infrared Spectrometry (FTIR)
%eveloped a device capable of retaining a suspended cells, as small as 6 um in diameter, while permitting
%luid transport with the purpose of studying cell membrane stability.
%ydraulics Laboratory (Dept. of B.B.E.) Summer 2007 – Fall 2008
%n charge of the research and development of a subterranean fluxmeter. The purpose of the device is to
%easure fluid flow while being minimally invasive. The fluid was also stored in a compartment for
%osterior extraction and chemical analysis. The fluxmeter was designed to be capable of sensing flow
%alues comparable to those recorded by commercial rain gauges.
%ourses and skills
%roficient in Pro/Engineer, ANSYS, C, MATLAB and SimHydraulics.
%ublications
%andeep Dhull, Dario Canelon, Apostolos Kottas, Justin Dancs, Andrew Carlson, and Nikolaos
%apanikolopoulos. "Aquapod: A Small Amphibious Robot with Sampling Capabilities". IEEE
%nternational Conference on Intelligent Robots and Systems (IROS) 2012.
%rett Hemes, Dario Canelon, Justin Dancs and Nikolaos Papanikolopoulos "Robotic Tumbling
%ocomotion: Overview, the Adelopod-T, and Step Climbing Analysis". IEEE International
%onference on Robotics and Automation (ICRA) 2011. 
  
\end{singlespace}